\documentclass[aspectratio=169,10pt]{beamer}
\makeatletter\def\input@path{{../../estilo/}{../estilo/}{estilo/}}\makeatother
\usepackage{devmobile}
\modulo[Flutter]{07}{Estado, Navegação e Animações}{setState, ciclo de vida, elevação de estado, ChangeNotifier, ciclo de vida do app, Navigator, go\_router, diálogos e animações}

\begin{document}
\capa

\begin{frame}{Objetivos do módulo}
  \begin{conceito}[Ao final deste módulo você será capaz de]
    \begin{itemize}
      \item Diferenciar estado efêmero de estado do app e usar \cod{setState} corretamente.
      \item Explicar o ciclo de vida de um \cod{State} e do próprio app.
      \item Elevar o estado e compartilhá-lo com \cod{ValueNotifier} e \cod{ChangeNotifier}.
      \item Navegar entre telas com \cod{Navigator} e com \textbf{go\_router} (rotas, parâmetros, abas e proteção por login).
      \item Mostrar diálogos, \emph{snackbars} e \emph{bottom sheets}.
      \item Criar animações implícitas, explícitas e transições com \cod{Hero}.
    \end{itemize}
  \end{conceito}
\end{frame}

\sumario

% =============================================================================
%  Fundamentos compartilhados --- Módulo 07: estado, ciclo de vida, navegação
%  e movimento
%  Incluído pelas trilhas Kotlin e Flutter com % =============================================================================
%  Fundamentos compartilhados --- Módulo 07: estado, ciclo de vida, navegação
%  e movimento
%  Incluído pelas trilhas Kotlin e Flutter com % =============================================================================
%  Fundamentos compartilhados --- Módulo 07: estado, ciclo de vida, navegação
%  e movimento
%  Incluído pelas trilhas Kotlin e Flutter com \input{../fundamentos/f07-estado-navegacao}
% =============================================================================
\section{Fundamentos: estado, ciclo de vida e navegação}

\begin{frame}{Estado: o que a tela precisa lembrar}
  \begin{columns}[T]
    \begin{column}{0.56\textwidth}
      \begin{definicao}[Estado]
        Todo valor que pode mudar com o tempo e influencia o que a tela mostra. A interface é uma função dele: \textbf{UI = f(estado)}.
      \end{definicao}
      \scriptsize
      \begin{tabularx}{\linewidth}{@{}l L l@{}}
        \toprule
        \cabecalho{Tipo} & \cabecalho{Exemplo no LabMonitor} & \cabecalho{Onde mora} \\
        \midrule
        Efêmero & painel expandido, rolagem, aba ativa & no componente \\
        Da tela & bancadas carregadas, filtro, erro & no dono da tela \\
        Do app & usuário logado, tema escolhido & em um serviço compartilhado \\
        Persistente & reservas, preferências & banco ou disco (Módulo 10) \\
        \bottomrule
      \end{tabularx}
    \end{column}
    \begin{column}{0.42\textwidth}
      \begin{definicao}[Estado derivado]
        Valor calculado a partir de outros, como o total de bancadas livres. Não se guarda: calcula-se, para nunca ficar desatualizado.
      \end{definicao}
      \begin{conceito}[Fonte única da verdade]
        Cada dado tem um único dono que pode alterá-lo; os demais apenas leem. Duas cópias editáveis do mesmo dado acabam divergindo.
      \end{conceito}
    \end{column}
  \end{columns}
\end{frame}

\begin{frame}{Fluxo unidirecional de dados}
  \begin{columns}[c]
    \begin{column}{0.46\textwidth}
      \centering
      \begin{tikzpicture}[every node/.style={font=\scriptsize}]
        \node[caixaambar,minimum width=2.4cm] (e) at (0,1.5) {\textbf{Estado}};
        \node[caixa,minimum width=2.4cm] (u) at (2,0) {\textbf{Interface}};
        \node[caixa,minimum width=2.4cm] (v) at (0,-1.5) {\textbf{Evento}\\ \tiny toque, digitação};
        \node[caixaescura,minimum width=2.4cm] (d) at (-2,0) {\textbf{Dono do estado}\\ \tiny aplica as regras};
        \draw[seta] (e) to[bend left=20] node[rotulo,right]{desce} (u);
        \draw[setaambar] (u) to[bend left=20] node[rotulo,right]{gera} (v);
        \draw[setaambar] (v) to[bend left=20] node[rotulo,left]{sobe} (d);
        \draw[seta] (d) to[bend left=20] node[rotulo,left]{novo valor} (e);
      \end{tikzpicture}
    \end{column}
    \begin{column}{0.52\textwidth}
      \begin{definicao}[Fluxo unidirecional (UDF)]
        O estado desce do dono até a interface; os eventos sobem da interface até o dono. \textbf{Só o dono altera o estado}.
      \end{definicao}
      \begin{conceito}[Vantagens]
        Previsível: cada mudança tem um único caminho. Fácil de depurar: registre eventos e estados. Fácil de testar: um evento entra, um estado sai.
      \end{conceito}
      \begin{historia}
        A Arquitetura Elm, o Flux do Facebook (2014) e o Redux (2015) popularizaram o modelo, hoje padrão no Compose, no Flutter e no SwiftUI.
      \end{historia}
    \end{column}
  \end{columns}
\end{frame}

\begin{frame}{A tela como máquina de estados}
  \begin{columns}[c]
    \begin{column}{0.5\textwidth}
      \centering
      \begin{tikzpicture}[every node/.style={font=\scriptsize}]
        \node[caixaescura,minimum width=2cm] (c) at (0,0) {Carregando};
        \node[caixa,minimum width=1.8cm,fill=verdeok!12,draw=verdeok] (s) at (3,1.1) {Sucesso};
        \node[caixa,minimum width=1.8cm] (v) at (3,0) {Vazio};
        \node[caixa,minimum width=1.8cm,fill=vermelhoerro!10,draw=vermelhoerro] (x) at (3,-1.1) {Erro};
        \draw[seta] (-1.6,0) node[rotulo,left]{abrir tela} -- (c);
        \draw[seta] (c) -- node[rotulo,above,sloped]{dados} (s.west);
        \draw[seta] (c) -- node[rotulo,above]{lista vazia} (v);
        \draw[seta] (c) -- node[rotulo,below,sloped]{falha} (x.west);
        \draw[setaambar,rounded corners] (x.south) -- ++(0,-0.4) -| node[rotulo,pos=0.25,below]{tentar de novo} (c.south);
        \draw[setaambar,rounded corners] (s.north) -- ++(0,0.4) -| node[rotulo,pos=0.25,above]{puxar para atualizar} (c.north);
      \end{tikzpicture}
    \end{column}
    \begin{column}{0.48\textwidth}
      \begin{definicao}[Máquina de estados finitos]
        Conjunto fechado de estados e de transições permitidas, disparadas por eventos. A tela está sempre em \textbf{exatamente um} estado.
      \end{definicao}
      \begin{atencao}[Estados impossíveis]
        Três booleanos (\cod{carregando}, \cod{erro}, \cod{vazio}) geram 8 combinações, a maioria sem sentido. Uma classe selada com quatro casos admite só os quatro válidos.
      \end{atencao}
      \begin{dica}
        Kotlin e Dart 3 têm classes seladas: o compilador obriga o \cod{when} ou o \cod{switch} a tratar todos os casos.
      \end{dica}
    \end{column}
  \end{columns}
\end{frame}

\begin{frame}{O que sobrevive a quê}
  \scriptsize
  \newcommand{\vsim}{\textcolor{verdeok}{\textbf{sim}}}
  \newcommand{\vnao}{\textcolor{vermelhoerro}{não}}
  \begin{tabularx}{\textwidth}{@{}L *{4}{>{\centering\arraybackslash}p{1.95cm}}@{}}
    \toprule
    \cabecalho{Onde o valor está} & \cabecalho{nova recomposição} & \cabecalho{rotação da tela} & \cabecalho{morte do processo} & \cabecalho{app fechado pelo usuário} \\
    \midrule
    variável local dentro da função de interface & \vnao & \vnao & \vnao & \vnao \\
    \cod{remember} (Compose) & \vsim & \vnao & \vnao & \vnao \\
    \cod{rememberSaveable} (Compose) & \vsim & \vsim & \vsim & \vnao \\
    \cod{ViewModel} (Android) & \vsim & \vsim & \vnao$^1$ & \vnao \\
    \cod{State} de um \cod{StatefulWidget} (Flutter) & \vsim & \vsim$^2$ & \vnao & \vnao \\
    \cod{RestorationMixin} (Flutter) & \vsim & \vsim & \vsim & \vnao \\
    banco de dados, arquivo, preferências (Módulo 10) & \vsim & \vsim & \vsim & \vsim \\
    \bottomrule
  \end{tabularx}\\[1pt]
  {\tiny $^1$ sim, com \cod{SavedStateHandle}. \quad $^2$ o Flutter não recria a tela ao girar.}
  \normalsize
  \begin{columns}[T]
    \begin{column}{0.49\textwidth}
      \begin{definicao}[Morte do processo]
        Em segundo plano, o sistema pode encerrar o app para liberar memória. Ao voltar, o usuário espera a tela como a deixou: o sistema a recria e devolve o pequeno estado salvo.
      \end{definicao}
    \end{column}
    \begin{column}{0.49\textwidth}
      \begin{dica}[Como testar]
        Ative ``Não manter atividades'' nas opções do desenvolvedor do Android, ou rode \cod{adb shell am kill <pacote>} com o app em segundo plano.
      \end{dica}
    \end{column}
  \end{columns}
\end{frame}

\begin{frame}{Ciclo de vida e prioridade dos processos}
  \begin{columns}[c]
    \begin{column}{0.44\textwidth}
      \centering
      \begin{tikzpicture}[node distance=1.5mm,every node/.style={font=\scriptsize}]
        \node[caixaescura,minimum width=3.6cm] (a) {primeiro plano (em uso)};
        \node[caixapetroleo,minimum width=3.6cm,below=of a] (b) {visível, mas encoberto};
        \node[caixa,minimum width=3.6cm,below=of b] (c) {serviço em execução};
        \node[caixa,minimum width=3.6cm,below=of c,fill=cinzaclaro] (d) {segundo plano (em cache)};
        \draw[seta] ([xshift=2mm]a.north east) -- node[rotulo,right,text width=1.4cm]{mais importante} ([xshift=2mm]d.south east);
        \node[rotulo,below=2mm of d,text width=3.8cm] {com pouca memória, o sistema encerra de baixo para cima};
      \end{tikzpicture}
    \end{column}
    \begin{column}{0.54\textwidth}
      \begin{definicao}[Ciclo de vida]
        Sequência de estados por que passam o app e cada tela, controlada pelo sistema: criado, visível, em primeiro plano, em segundo plano, destruído. Cada transição é avisada por \emph{callbacks}.
      \end{definicao}
      \begin{conceito}[O que fazer nas transições]
        Ao sair do primeiro plano: pare câmera, sensores e animações e salve rascunhos. Ao voltar: retome e atualize os dados.
      \end{conceito}
      \begin{pordentro}[Por que o sistema encerra apps]
        A memória do celular é limitada, e o sistema prefere encerrar apps em segundo plano a deixar o aparelho lento: no Android, o \emph{low memory killer}; no iOS, o \emph{jetsam}.
      \end{pordentro}
    \end{column}
  \end{columns}
\end{frame}

\begin{frame}{Navegação: grafo de destinos e pilha de retorno}
  \begin{columns}[c]
    \begin{column}{0.48\textwidth}
      \centering
      \begin{tikzpicture}[node distance=4mm and 3mm,every node/.style={font=\tiny}]
        \node[caixa,minimum width=1.4cm,font=\tiny] (l) {Login};
        \node[caixaescura,minimum width=1.4cm,font=\tiny,right=of l] (p) {Painel};
        \node[caixa,minimum width=1.4cm,font=\tiny,below=of p] (e) {Equipamentos};
        \node[caixa,minimum width=1.4cm,font=\tiny,right=of p] (a) {Ajustes};
        \node[caixa,minimum width=1.4cm,font=\tiny,below=of e] (d) {Detalhe};
        \node[caixaambar,minimum width=1.4cm,font=\tiny,below=of d] (r) {Reserva};
        \draw[seta] (l) -- (p); \draw[seta] (p) -- (e); \draw[seta] (p) -- (a); \draw[seta] (e) -- (d); \draw[seta] (d) -- (r);
        \node[rotulo,below=1mm of r] {grafo};
        \node[caixa,minimum width=1.9cm,font=\tiny] (s1) at (4.6,-2.6) {Painel};
        \node[caixa,minimum width=1.9cm,font=\tiny,above=0.5mm of s1] (s2) {Equipamentos};
        \node[caixaambarforte,minimum width=1.9cm,font=\tiny,above=0.5mm of s2] (s3) {Detalhe (topo)};
        \node[pontilhado,fit=(s1)(s2)(s3)] (pil) {};
        \node[rotulo,below=1mm of pil] {pilha};
      \end{tikzpicture}
    \end{column}
    \begin{column}{0.5\textwidth}
      \begin{definicao}[Grafo e pilha]
        Destinos (telas) são \textbf{nós}; as ligações permitidas, \textbf{arestas}. A pilha de retorno é LIFO: navegar empilha, voltar desempilha.
      \end{definicao}
      \begin{conceito}[Princípios]
        Destino inicial fixo; o voltar nunca surpreende; cada aba da barra inferior tem a própria pilha.
      \end{conceito}
      \begin{atencao}[Voltar preditivo]
        Com \cod{targetSdk} 36 no Android 16, o gesto mostra a tela anterior antes de o dedo soltar e \cod{onBackPressed} não é mais chamado: use \cod{BackHandler} ou \cod{PopScope}.
      \end{atencao}
    \end{column}
  \end{columns}
\end{frame}

\begin{frame}{Navegação declarativa e links profundos}
  \begin{columns}[T]
    \begin{column}{0.5\textwidth}
      \begin{definicao}[Imperativa \emph{versus} declarativa]
        Imperativa: comandos \cod{push} e \cod{pop}. Declarativa: a pilha é um \textbf{estado} (uma lista de destinos) e o que aparece é função dela, como na interface (Navigation 3, go\_router).
      \end{definicao}
      \begin{definicao}[Link profundo (\emph{deep link})]
        Endereço que abre o app direto em um destino, como \cod{https://labmonitor.app/reserva/42}. O app monta uma pilha sintética (Painel, Reserva) para o voltar fazer sentido.
      \end{definicao}
      \begin{atencao}
        Parâmetros de um link vêm de fora: valide-os (a reserva existe? este usuário pode vê-la?).
      \end{atencao}
    \end{column}
    \begin{column}{0.48\textwidth}
      \centering
      \begin{tikzpicture}[node distance=3mm,every node/.style={font=\tiny}]
        \node[caixa,minimum width=3.6cm,font=\tiny] (o) {e-mail, notificação ou QR Code};
        \node[caixapetroleo,minimum width=3.6cm,font=\tiny,below=of o] (s) {sistema confere o domínio};
        \node[caixaescura,minimum width=3.6cm,font=\tiny,below=of s] (r) {roteador do app interpreta};
        \node[caixaambar,minimum width=3.6cm,font=\tiny,below=of r] (p) {pilha: Painel $\rightarrow$ Reserva 42};
        \draw[seta] (o) -- (s); \draw[seta] (s) -- (r); \draw[seta] (r) -- (p);
      \end{tikzpicture}
      \raggedright
      \begin{pordentro}[Links verificados]
        \emph{App Links} (Android, arquivo \cod{assetlinks.json}) e \emph{Universal Links} (iOS, \cod{apple-app-site-association}) provam que o site e o app têm o mesmo dono.
      \end{pordentro}
    \end{column}
  \end{columns}
\end{frame}

\begin{frame}{Movimento: interpolação, curvas e molas}
  \begin{columns}[T]
    \begin{column}{0.44\textwidth}
      \centering
      \begin{tikzpicture}[x=3.4cm,y=1.9cm,baseline=(current bounding box.north)]
        \draw[->,cinzamedio] (0,0) -- (1.08,0) node[rotulo,right]{t};
        \draw[->,cinzamedio] (0,0) -- (0,1.35) node[rotulo,above]{progresso};
        \draw[cinzaclaro,dashed] (0,1) -- (1,1);
        \draw[cinzamedio,line width=0.8pt] plot[domain=0:1,samples=2] (\x,\x);
        \draw[azulpetroleo,line width=1pt] plot[domain=0:1,samples=40] (\x,{1-(1-\x)^3});
        \draw[ambar,line width=1pt] plot[domain=0:1,samples=80] (\x,{1-exp(-5*\x)*cos(deg(10*\x))});
        \node[rotulo,text=cinzamedio] at (0.82,0.62) {linear};
        \node[rotulo,text=azulpetroleo] at (0.3,0.88) {desacelera};
        \node[rotulo,text=ambar] at (0.42,1.3) {mola (passa e volta)};
      \end{tikzpicture}
      \raggedright
      \begin{definicao}[Interpolação]
        A cada quadro: \cod{valor = início + (fim - início)} $\times$ \cod{curva(t)}, com \cod{t} indo de 0 a 1 ao longo da animação.
      \end{definicao}
    \end{column}
    \begin{column}{0.54\textwidth}
      \begin{conceito}[Curvas e molas]
        Curvas de aceleração (\emph{easing}) imitam objetos reais: desaceleram ao chegar. \textbf{Molas} são definidas por rigidez e amortecimento, não por duração, e podem ser interrompidas sem saltos. O Material 3 Expressive (2025) adotou molas como base.
      \end{conceito}
      \begin{dica}[Movimento com propósito]
        Anime para orientar: de onde veio a tela, o que mudou. Transições de 150 a 400\,ms mantêm o app ágil.
      \end{dica}
      \begin{atencao}[Acessibilidade]
        Respeite a opção do sistema de remover ou reduzir animações: movimento excessivo causa desconforto em algumas pessoas.
      \end{atencao}
    \end{column}
  \end{columns}
\end{frame}

% =============================================================================
\section{Fundamentos: estado, ciclo de vida e navegação}

\begin{frame}{Estado: o que a tela precisa lembrar}
  \begin{columns}[T]
    \begin{column}{0.56\textwidth}
      \begin{definicao}[Estado]
        Todo valor que pode mudar com o tempo e influencia o que a tela mostra. A interface é uma função dele: \textbf{UI = f(estado)}.
      \end{definicao}
      \scriptsize
      \begin{tabularx}{\linewidth}{@{}l L l@{}}
        \toprule
        \cabecalho{Tipo} & \cabecalho{Exemplo no LabMonitor} & \cabecalho{Onde mora} \\
        \midrule
        Efêmero & painel expandido, rolagem, aba ativa & no componente \\
        Da tela & bancadas carregadas, filtro, erro & no dono da tela \\
        Do app & usuário logado, tema escolhido & em um serviço compartilhado \\
        Persistente & reservas, preferências & banco ou disco (Módulo 10) \\
        \bottomrule
      \end{tabularx}
    \end{column}
    \begin{column}{0.42\textwidth}
      \begin{definicao}[Estado derivado]
        Valor calculado a partir de outros, como o total de bancadas livres. Não se guarda: calcula-se, para nunca ficar desatualizado.
      \end{definicao}
      \begin{conceito}[Fonte única da verdade]
        Cada dado tem um único dono que pode alterá-lo; os demais apenas leem. Duas cópias editáveis do mesmo dado acabam divergindo.
      \end{conceito}
    \end{column}
  \end{columns}
\end{frame}

\begin{frame}{Fluxo unidirecional de dados}
  \begin{columns}[c]
    \begin{column}{0.46\textwidth}
      \centering
      \begin{tikzpicture}[every node/.style={font=\scriptsize}]
        \node[caixaambar,minimum width=2.4cm] (e) at (0,1.5) {\textbf{Estado}};
        \node[caixa,minimum width=2.4cm] (u) at (2,0) {\textbf{Interface}};
        \node[caixa,minimum width=2.4cm] (v) at (0,-1.5) {\textbf{Evento}\\ \tiny toque, digitação};
        \node[caixaescura,minimum width=2.4cm] (d) at (-2,0) {\textbf{Dono do estado}\\ \tiny aplica as regras};
        \draw[seta] (e) to[bend left=20] node[rotulo,right]{desce} (u);
        \draw[setaambar] (u) to[bend left=20] node[rotulo,right]{gera} (v);
        \draw[setaambar] (v) to[bend left=20] node[rotulo,left]{sobe} (d);
        \draw[seta] (d) to[bend left=20] node[rotulo,left]{novo valor} (e);
      \end{tikzpicture}
    \end{column}
    \begin{column}{0.52\textwidth}
      \begin{definicao}[Fluxo unidirecional (UDF)]
        O estado desce do dono até a interface; os eventos sobem da interface até o dono. \textbf{Só o dono altera o estado}.
      \end{definicao}
      \begin{conceito}[Vantagens]
        Previsível: cada mudança tem um único caminho. Fácil de depurar: registre eventos e estados. Fácil de testar: um evento entra, um estado sai.
      \end{conceito}
      \begin{historia}
        A Arquitetura Elm, o Flux do Facebook (2014) e o Redux (2015) popularizaram o modelo, hoje padrão no Compose, no Flutter e no SwiftUI.
      \end{historia}
    \end{column}
  \end{columns}
\end{frame}

\begin{frame}{A tela como máquina de estados}
  \begin{columns}[c]
    \begin{column}{0.5\textwidth}
      \centering
      \begin{tikzpicture}[every node/.style={font=\scriptsize}]
        \node[caixaescura,minimum width=2cm] (c) at (0,0) {Carregando};
        \node[caixa,minimum width=1.8cm,fill=verdeok!12,draw=verdeok] (s) at (3,1.1) {Sucesso};
        \node[caixa,minimum width=1.8cm] (v) at (3,0) {Vazio};
        \node[caixa,minimum width=1.8cm,fill=vermelhoerro!10,draw=vermelhoerro] (x) at (3,-1.1) {Erro};
        \draw[seta] (-1.6,0) node[rotulo,left]{abrir tela} -- (c);
        \draw[seta] (c) -- node[rotulo,above,sloped]{dados} (s.west);
        \draw[seta] (c) -- node[rotulo,above]{lista vazia} (v);
        \draw[seta] (c) -- node[rotulo,below,sloped]{falha} (x.west);
        \draw[setaambar,rounded corners] (x.south) -- ++(0,-0.4) -| node[rotulo,pos=0.25,below]{tentar de novo} (c.south);
        \draw[setaambar,rounded corners] (s.north) -- ++(0,0.4) -| node[rotulo,pos=0.25,above]{puxar para atualizar} (c.north);
      \end{tikzpicture}
    \end{column}
    \begin{column}{0.48\textwidth}
      \begin{definicao}[Máquina de estados finitos]
        Conjunto fechado de estados e de transições permitidas, disparadas por eventos. A tela está sempre em \textbf{exatamente um} estado.
      \end{definicao}
      \begin{atencao}[Estados impossíveis]
        Três booleanos (\cod{carregando}, \cod{erro}, \cod{vazio}) geram 8 combinações, a maioria sem sentido. Uma classe selada com quatro casos admite só os quatro válidos.
      \end{atencao}
      \begin{dica}
        Kotlin e Dart 3 têm classes seladas: o compilador obriga o \cod{when} ou o \cod{switch} a tratar todos os casos.
      \end{dica}
    \end{column}
  \end{columns}
\end{frame}

\begin{frame}{O que sobrevive a quê}
  \scriptsize
  \newcommand{\vsim}{\textcolor{verdeok}{\textbf{sim}}}
  \newcommand{\vnao}{\textcolor{vermelhoerro}{não}}
  \begin{tabularx}{\textwidth}{@{}L *{4}{>{\centering\arraybackslash}p{1.95cm}}@{}}
    \toprule
    \cabecalho{Onde o valor está} & \cabecalho{nova recomposição} & \cabecalho{rotação da tela} & \cabecalho{morte do processo} & \cabecalho{app fechado pelo usuário} \\
    \midrule
    variável local dentro da função de interface & \vnao & \vnao & \vnao & \vnao \\
    \cod{remember} (Compose) & \vsim & \vnao & \vnao & \vnao \\
    \cod{rememberSaveable} (Compose) & \vsim & \vsim & \vsim & \vnao \\
    \cod{ViewModel} (Android) & \vsim & \vsim & \vnao$^1$ & \vnao \\
    \cod{State} de um \cod{StatefulWidget} (Flutter) & \vsim & \vsim$^2$ & \vnao & \vnao \\
    \cod{RestorationMixin} (Flutter) & \vsim & \vsim & \vsim & \vnao \\
    banco de dados, arquivo, preferências (Módulo 10) & \vsim & \vsim & \vsim & \vsim \\
    \bottomrule
  \end{tabularx}\\[1pt]
  {\tiny $^1$ sim, com \cod{SavedStateHandle}. \quad $^2$ o Flutter não recria a tela ao girar.}
  \normalsize
  \begin{columns}[T]
    \begin{column}{0.49\textwidth}
      \begin{definicao}[Morte do processo]
        Em segundo plano, o sistema pode encerrar o app para liberar memória. Ao voltar, o usuário espera a tela como a deixou: o sistema a recria e devolve o pequeno estado salvo.
      \end{definicao}
    \end{column}
    \begin{column}{0.49\textwidth}
      \begin{dica}[Como testar]
        Ative ``Não manter atividades'' nas opções do desenvolvedor do Android, ou rode \cod{adb shell am kill <pacote>} com o app em segundo plano.
      \end{dica}
    \end{column}
  \end{columns}
\end{frame}

\begin{frame}{Ciclo de vida e prioridade dos processos}
  \begin{columns}[c]
    \begin{column}{0.44\textwidth}
      \centering
      \begin{tikzpicture}[node distance=1.5mm,every node/.style={font=\scriptsize}]
        \node[caixaescura,minimum width=3.6cm] (a) {primeiro plano (em uso)};
        \node[caixapetroleo,minimum width=3.6cm,below=of a] (b) {visível, mas encoberto};
        \node[caixa,minimum width=3.6cm,below=of b] (c) {serviço em execução};
        \node[caixa,minimum width=3.6cm,below=of c,fill=cinzaclaro] (d) {segundo plano (em cache)};
        \draw[seta] ([xshift=2mm]a.north east) -- node[rotulo,right,text width=1.4cm]{mais importante} ([xshift=2mm]d.south east);
        \node[rotulo,below=2mm of d,text width=3.8cm] {com pouca memória, o sistema encerra de baixo para cima};
      \end{tikzpicture}
    \end{column}
    \begin{column}{0.54\textwidth}
      \begin{definicao}[Ciclo de vida]
        Sequência de estados por que passam o app e cada tela, controlada pelo sistema: criado, visível, em primeiro plano, em segundo plano, destruído. Cada transição é avisada por \emph{callbacks}.
      \end{definicao}
      \begin{conceito}[O que fazer nas transições]
        Ao sair do primeiro plano: pare câmera, sensores e animações e salve rascunhos. Ao voltar: retome e atualize os dados.
      \end{conceito}
      \begin{pordentro}[Por que o sistema encerra apps]
        A memória do celular é limitada, e o sistema prefere encerrar apps em segundo plano a deixar o aparelho lento: no Android, o \emph{low memory killer}; no iOS, o \emph{jetsam}.
      \end{pordentro}
    \end{column}
  \end{columns}
\end{frame}

\begin{frame}{Navegação: grafo de destinos e pilha de retorno}
  \begin{columns}[c]
    \begin{column}{0.48\textwidth}
      \centering
      \begin{tikzpicture}[node distance=4mm and 3mm,every node/.style={font=\tiny}]
        \node[caixa,minimum width=1.4cm,font=\tiny] (l) {Login};
        \node[caixaescura,minimum width=1.4cm,font=\tiny,right=of l] (p) {Painel};
        \node[caixa,minimum width=1.4cm,font=\tiny,below=of p] (e) {Equipamentos};
        \node[caixa,minimum width=1.4cm,font=\tiny,right=of p] (a) {Ajustes};
        \node[caixa,minimum width=1.4cm,font=\tiny,below=of e] (d) {Detalhe};
        \node[caixaambar,minimum width=1.4cm,font=\tiny,below=of d] (r) {Reserva};
        \draw[seta] (l) -- (p); \draw[seta] (p) -- (e); \draw[seta] (p) -- (a); \draw[seta] (e) -- (d); \draw[seta] (d) -- (r);
        \node[rotulo,below=1mm of r] {grafo};
        \node[caixa,minimum width=1.9cm,font=\tiny] (s1) at (4.6,-2.6) {Painel};
        \node[caixa,minimum width=1.9cm,font=\tiny,above=0.5mm of s1] (s2) {Equipamentos};
        \node[caixaambarforte,minimum width=1.9cm,font=\tiny,above=0.5mm of s2] (s3) {Detalhe (topo)};
        \node[pontilhado,fit=(s1)(s2)(s3)] (pil) {};
        \node[rotulo,below=1mm of pil] {pilha};
      \end{tikzpicture}
    \end{column}
    \begin{column}{0.5\textwidth}
      \begin{definicao}[Grafo e pilha]
        Destinos (telas) são \textbf{nós}; as ligações permitidas, \textbf{arestas}. A pilha de retorno é LIFO: navegar empilha, voltar desempilha.
      \end{definicao}
      \begin{conceito}[Princípios]
        Destino inicial fixo; o voltar nunca surpreende; cada aba da barra inferior tem a própria pilha.
      \end{conceito}
      \begin{atencao}[Voltar preditivo]
        Com \cod{targetSdk} 36 no Android 16, o gesto mostra a tela anterior antes de o dedo soltar e \cod{onBackPressed} não é mais chamado: use \cod{BackHandler} ou \cod{PopScope}.
      \end{atencao}
    \end{column}
  \end{columns}
\end{frame}

\begin{frame}{Navegação declarativa e links profundos}
  \begin{columns}[T]
    \begin{column}{0.5\textwidth}
      \begin{definicao}[Imperativa \emph{versus} declarativa]
        Imperativa: comandos \cod{push} e \cod{pop}. Declarativa: a pilha é um \textbf{estado} (uma lista de destinos) e o que aparece é função dela, como na interface (Navigation 3, go\_router).
      \end{definicao}
      \begin{definicao}[Link profundo (\emph{deep link})]
        Endereço que abre o app direto em um destino, como \cod{https://labmonitor.app/reserva/42}. O app monta uma pilha sintética (Painel, Reserva) para o voltar fazer sentido.
      \end{definicao}
      \begin{atencao}
        Parâmetros de um link vêm de fora: valide-os (a reserva existe? este usuário pode vê-la?).
      \end{atencao}
    \end{column}
    \begin{column}{0.48\textwidth}
      \centering
      \begin{tikzpicture}[node distance=3mm,every node/.style={font=\tiny}]
        \node[caixa,minimum width=3.6cm,font=\tiny] (o) {e-mail, notificação ou QR Code};
        \node[caixapetroleo,minimum width=3.6cm,font=\tiny,below=of o] (s) {sistema confere o domínio};
        \node[caixaescura,minimum width=3.6cm,font=\tiny,below=of s] (r) {roteador do app interpreta};
        \node[caixaambar,minimum width=3.6cm,font=\tiny,below=of r] (p) {pilha: Painel $\rightarrow$ Reserva 42};
        \draw[seta] (o) -- (s); \draw[seta] (s) -- (r); \draw[seta] (r) -- (p);
      \end{tikzpicture}
      \raggedright
      \begin{pordentro}[Links verificados]
        \emph{App Links} (Android, arquivo \cod{assetlinks.json}) e \emph{Universal Links} (iOS, \cod{apple-app-site-association}) provam que o site e o app têm o mesmo dono.
      \end{pordentro}
    \end{column}
  \end{columns}
\end{frame}

\begin{frame}{Movimento: interpolação, curvas e molas}
  \begin{columns}[T]
    \begin{column}{0.44\textwidth}
      \centering
      \begin{tikzpicture}[x=3.4cm,y=1.9cm,baseline=(current bounding box.north)]
        \draw[->,cinzamedio] (0,0) -- (1.08,0) node[rotulo,right]{t};
        \draw[->,cinzamedio] (0,0) -- (0,1.35) node[rotulo,above]{progresso};
        \draw[cinzaclaro,dashed] (0,1) -- (1,1);
        \draw[cinzamedio,line width=0.8pt] plot[domain=0:1,samples=2] (\x,\x);
        \draw[azulpetroleo,line width=1pt] plot[domain=0:1,samples=40] (\x,{1-(1-\x)^3});
        \draw[ambar,line width=1pt] plot[domain=0:1,samples=80] (\x,{1-exp(-5*\x)*cos(deg(10*\x))});
        \node[rotulo,text=cinzamedio] at (0.82,0.62) {linear};
        \node[rotulo,text=azulpetroleo] at (0.3,0.88) {desacelera};
        \node[rotulo,text=ambar] at (0.42,1.3) {mola (passa e volta)};
      \end{tikzpicture}
      \raggedright
      \begin{definicao}[Interpolação]
        A cada quadro: \cod{valor = início + (fim - início)} $\times$ \cod{curva(t)}, com \cod{t} indo de 0 a 1 ao longo da animação.
      \end{definicao}
    \end{column}
    \begin{column}{0.54\textwidth}
      \begin{conceito}[Curvas e molas]
        Curvas de aceleração (\emph{easing}) imitam objetos reais: desaceleram ao chegar. \textbf{Molas} são definidas por rigidez e amortecimento, não por duração, e podem ser interrompidas sem saltos. O Material 3 Expressive (2025) adotou molas como base.
      \end{conceito}
      \begin{dica}[Movimento com propósito]
        Anime para orientar: de onde veio a tela, o que mudou. Transições de 150 a 400\,ms mantêm o app ágil.
      \end{dica}
      \begin{atencao}[Acessibilidade]
        Respeite a opção do sistema de remover ou reduzir animações: movimento excessivo causa desconforto em algumas pessoas.
      \end{atencao}
    \end{column}
  \end{columns}
\end{frame}

% =============================================================================
\section{Fundamentos: estado, ciclo de vida e navegação}

\begin{frame}{Estado: o que a tela precisa lembrar}
  \begin{columns}[T]
    \begin{column}{0.56\textwidth}
      \begin{definicao}[Estado]
        Todo valor que pode mudar com o tempo e influencia o que a tela mostra. A interface é uma função dele: \textbf{UI = f(estado)}.
      \end{definicao}
      \scriptsize
      \begin{tabularx}{\linewidth}{@{}l L l@{}}
        \toprule
        \cabecalho{Tipo} & \cabecalho{Exemplo no LabMonitor} & \cabecalho{Onde mora} \\
        \midrule
        Efêmero & painel expandido, rolagem, aba ativa & no componente \\
        Da tela & bancadas carregadas, filtro, erro & no dono da tela \\
        Do app & usuário logado, tema escolhido & em um serviço compartilhado \\
        Persistente & reservas, preferências & banco ou disco (Módulo 10) \\
        \bottomrule
      \end{tabularx}
    \end{column}
    \begin{column}{0.42\textwidth}
      \begin{definicao}[Estado derivado]
        Valor calculado a partir de outros, como o total de bancadas livres. Não se guarda: calcula-se, para nunca ficar desatualizado.
      \end{definicao}
      \begin{conceito}[Fonte única da verdade]
        Cada dado tem um único dono que pode alterá-lo; os demais apenas leem. Duas cópias editáveis do mesmo dado acabam divergindo.
      \end{conceito}
    \end{column}
  \end{columns}
\end{frame}

\begin{frame}{Fluxo unidirecional de dados}
  \begin{columns}[c]
    \begin{column}{0.46\textwidth}
      \centering
      \begin{tikzpicture}[every node/.style={font=\scriptsize}]
        \node[caixaambar,minimum width=2.4cm] (e) at (0,1.5) {\textbf{Estado}};
        \node[caixa,minimum width=2.4cm] (u) at (2,0) {\textbf{Interface}};
        \node[caixa,minimum width=2.4cm] (v) at (0,-1.5) {\textbf{Evento}\\ \tiny toque, digitação};
        \node[caixaescura,minimum width=2.4cm] (d) at (-2,0) {\textbf{Dono do estado}\\ \tiny aplica as regras};
        \draw[seta] (e) to[bend left=20] node[rotulo,right]{desce} (u);
        \draw[setaambar] (u) to[bend left=20] node[rotulo,right]{gera} (v);
        \draw[setaambar] (v) to[bend left=20] node[rotulo,left]{sobe} (d);
        \draw[seta] (d) to[bend left=20] node[rotulo,left]{novo valor} (e);
      \end{tikzpicture}
    \end{column}
    \begin{column}{0.52\textwidth}
      \begin{definicao}[Fluxo unidirecional (UDF)]
        O estado desce do dono até a interface; os eventos sobem da interface até o dono. \textbf{Só o dono altera o estado}.
      \end{definicao}
      \begin{conceito}[Vantagens]
        Previsível: cada mudança tem um único caminho. Fácil de depurar: registre eventos e estados. Fácil de testar: um evento entra, um estado sai.
      \end{conceito}
      \begin{historia}
        A Arquitetura Elm, o Flux do Facebook (2014) e o Redux (2015) popularizaram o modelo, hoje padrão no Compose, no Flutter e no SwiftUI.
      \end{historia}
    \end{column}
  \end{columns}
\end{frame}

\begin{frame}{A tela como máquina de estados}
  \begin{columns}[c]
    \begin{column}{0.5\textwidth}
      \centering
      \begin{tikzpicture}[every node/.style={font=\scriptsize}]
        \node[caixaescura,minimum width=2cm] (c) at (0,0) {Carregando};
        \node[caixa,minimum width=1.8cm,fill=verdeok!12,draw=verdeok] (s) at (3,1.1) {Sucesso};
        \node[caixa,minimum width=1.8cm] (v) at (3,0) {Vazio};
        \node[caixa,minimum width=1.8cm,fill=vermelhoerro!10,draw=vermelhoerro] (x) at (3,-1.1) {Erro};
        \draw[seta] (-1.6,0) node[rotulo,left]{abrir tela} -- (c);
        \draw[seta] (c) -- node[rotulo,above,sloped]{dados} (s.west);
        \draw[seta] (c) -- node[rotulo,above]{lista vazia} (v);
        \draw[seta] (c) -- node[rotulo,below,sloped]{falha} (x.west);
        \draw[setaambar,rounded corners] (x.south) -- ++(0,-0.4) -| node[rotulo,pos=0.25,below]{tentar de novo} (c.south);
        \draw[setaambar,rounded corners] (s.north) -- ++(0,0.4) -| node[rotulo,pos=0.25,above]{puxar para atualizar} (c.north);
      \end{tikzpicture}
    \end{column}
    \begin{column}{0.48\textwidth}
      \begin{definicao}[Máquina de estados finitos]
        Conjunto fechado de estados e de transições permitidas, disparadas por eventos. A tela está sempre em \textbf{exatamente um} estado.
      \end{definicao}
      \begin{atencao}[Estados impossíveis]
        Três booleanos (\cod{carregando}, \cod{erro}, \cod{vazio}) geram 8 combinações, a maioria sem sentido. Uma classe selada com quatro casos admite só os quatro válidos.
      \end{atencao}
      \begin{dica}
        Kotlin e Dart 3 têm classes seladas: o compilador obriga o \cod{when} ou o \cod{switch} a tratar todos os casos.
      \end{dica}
    \end{column}
  \end{columns}
\end{frame}

\begin{frame}{O que sobrevive a quê}
  \scriptsize
  \newcommand{\vsim}{\textcolor{verdeok}{\textbf{sim}}}
  \newcommand{\vnao}{\textcolor{vermelhoerro}{não}}
  \begin{tabularx}{\textwidth}{@{}L *{4}{>{\centering\arraybackslash}p{1.95cm}}@{}}
    \toprule
    \cabecalho{Onde o valor está} & \cabecalho{nova recomposição} & \cabecalho{rotação da tela} & \cabecalho{morte do processo} & \cabecalho{app fechado pelo usuário} \\
    \midrule
    variável local dentro da função de interface & \vnao & \vnao & \vnao & \vnao \\
    \cod{remember} (Compose) & \vsim & \vnao & \vnao & \vnao \\
    \cod{rememberSaveable} (Compose) & \vsim & \vsim & \vsim & \vnao \\
    \cod{ViewModel} (Android) & \vsim & \vsim & \vnao$^1$ & \vnao \\
    \cod{State} de um \cod{StatefulWidget} (Flutter) & \vsim & \vsim$^2$ & \vnao & \vnao \\
    \cod{RestorationMixin} (Flutter) & \vsim & \vsim & \vsim & \vnao \\
    banco de dados, arquivo, preferências (Módulo 10) & \vsim & \vsim & \vsim & \vsim \\
    \bottomrule
  \end{tabularx}\\[1pt]
  {\tiny $^1$ sim, com \cod{SavedStateHandle}. \quad $^2$ o Flutter não recria a tela ao girar.}
  \normalsize
  \begin{columns}[T]
    \begin{column}{0.49\textwidth}
      \begin{definicao}[Morte do processo]
        Em segundo plano, o sistema pode encerrar o app para liberar memória. Ao voltar, o usuário espera a tela como a deixou: o sistema a recria e devolve o pequeno estado salvo.
      \end{definicao}
    \end{column}
    \begin{column}{0.49\textwidth}
      \begin{dica}[Como testar]
        Ative ``Não manter atividades'' nas opções do desenvolvedor do Android, ou rode \cod{adb shell am kill <pacote>} com o app em segundo plano.
      \end{dica}
    \end{column}
  \end{columns}
\end{frame}

\begin{frame}{Ciclo de vida e prioridade dos processos}
  \begin{columns}[c]
    \begin{column}{0.44\textwidth}
      \centering
      \begin{tikzpicture}[node distance=1.5mm,every node/.style={font=\scriptsize}]
        \node[caixaescura,minimum width=3.6cm] (a) {primeiro plano (em uso)};
        \node[caixapetroleo,minimum width=3.6cm,below=of a] (b) {visível, mas encoberto};
        \node[caixa,minimum width=3.6cm,below=of b] (c) {serviço em execução};
        \node[caixa,minimum width=3.6cm,below=of c,fill=cinzaclaro] (d) {segundo plano (em cache)};
        \draw[seta] ([xshift=2mm]a.north east) -- node[rotulo,right,text width=1.4cm]{mais importante} ([xshift=2mm]d.south east);
        \node[rotulo,below=2mm of d,text width=3.8cm] {com pouca memória, o sistema encerra de baixo para cima};
      \end{tikzpicture}
    \end{column}
    \begin{column}{0.54\textwidth}
      \begin{definicao}[Ciclo de vida]
        Sequência de estados por que passam o app e cada tela, controlada pelo sistema: criado, visível, em primeiro plano, em segundo plano, destruído. Cada transição é avisada por \emph{callbacks}.
      \end{definicao}
      \begin{conceito}[O que fazer nas transições]
        Ao sair do primeiro plano: pare câmera, sensores e animações e salve rascunhos. Ao voltar: retome e atualize os dados.
      \end{conceito}
      \begin{pordentro}[Por que o sistema encerra apps]
        A memória do celular é limitada, e o sistema prefere encerrar apps em segundo plano a deixar o aparelho lento: no Android, o \emph{low memory killer}; no iOS, o \emph{jetsam}.
      \end{pordentro}
    \end{column}
  \end{columns}
\end{frame}

\begin{frame}{Navegação: grafo de destinos e pilha de retorno}
  \begin{columns}[c]
    \begin{column}{0.48\textwidth}
      \centering
      \begin{tikzpicture}[node distance=4mm and 3mm,every node/.style={font=\tiny}]
        \node[caixa,minimum width=1.4cm,font=\tiny] (l) {Login};
        \node[caixaescura,minimum width=1.4cm,font=\tiny,right=of l] (p) {Painel};
        \node[caixa,minimum width=1.4cm,font=\tiny,below=of p] (e) {Equipamentos};
        \node[caixa,minimum width=1.4cm,font=\tiny,right=of p] (a) {Ajustes};
        \node[caixa,minimum width=1.4cm,font=\tiny,below=of e] (d) {Detalhe};
        \node[caixaambar,minimum width=1.4cm,font=\tiny,below=of d] (r) {Reserva};
        \draw[seta] (l) -- (p); \draw[seta] (p) -- (e); \draw[seta] (p) -- (a); \draw[seta] (e) -- (d); \draw[seta] (d) -- (r);
        \node[rotulo,below=1mm of r] {grafo};
        \node[caixa,minimum width=1.9cm,font=\tiny] (s1) at (4.6,-2.6) {Painel};
        \node[caixa,minimum width=1.9cm,font=\tiny,above=0.5mm of s1] (s2) {Equipamentos};
        \node[caixaambarforte,minimum width=1.9cm,font=\tiny,above=0.5mm of s2] (s3) {Detalhe (topo)};
        \node[pontilhado,fit=(s1)(s2)(s3)] (pil) {};
        \node[rotulo,below=1mm of pil] {pilha};
      \end{tikzpicture}
    \end{column}
    \begin{column}{0.5\textwidth}
      \begin{definicao}[Grafo e pilha]
        Destinos (telas) são \textbf{nós}; as ligações permitidas, \textbf{arestas}. A pilha de retorno é LIFO: navegar empilha, voltar desempilha.
      \end{definicao}
      \begin{conceito}[Princípios]
        Destino inicial fixo; o voltar nunca surpreende; cada aba da barra inferior tem a própria pilha.
      \end{conceito}
      \begin{atencao}[Voltar preditivo]
        Com \cod{targetSdk} 36 no Android 16, o gesto mostra a tela anterior antes de o dedo soltar e \cod{onBackPressed} não é mais chamado: use \cod{BackHandler} ou \cod{PopScope}.
      \end{atencao}
    \end{column}
  \end{columns}
\end{frame}

\begin{frame}{Navegação declarativa e links profundos}
  \begin{columns}[T]
    \begin{column}{0.5\textwidth}
      \begin{definicao}[Imperativa \emph{versus} declarativa]
        Imperativa: comandos \cod{push} e \cod{pop}. Declarativa: a pilha é um \textbf{estado} (uma lista de destinos) e o que aparece é função dela, como na interface (Navigation 3, go\_router).
      \end{definicao}
      \begin{definicao}[Link profundo (\emph{deep link})]
        Endereço que abre o app direto em um destino, como \cod{https://labmonitor.app/reserva/42}. O app monta uma pilha sintética (Painel, Reserva) para o voltar fazer sentido.
      \end{definicao}
      \begin{atencao}
        Parâmetros de um link vêm de fora: valide-os (a reserva existe? este usuário pode vê-la?).
      \end{atencao}
    \end{column}
    \begin{column}{0.48\textwidth}
      \centering
      \begin{tikzpicture}[node distance=3mm,every node/.style={font=\tiny}]
        \node[caixa,minimum width=3.6cm,font=\tiny] (o) {e-mail, notificação ou QR Code};
        \node[caixapetroleo,minimum width=3.6cm,font=\tiny,below=of o] (s) {sistema confere o domínio};
        \node[caixaescura,minimum width=3.6cm,font=\tiny,below=of s] (r) {roteador do app interpreta};
        \node[caixaambar,minimum width=3.6cm,font=\tiny,below=of r] (p) {pilha: Painel $\rightarrow$ Reserva 42};
        \draw[seta] (o) -- (s); \draw[seta] (s) -- (r); \draw[seta] (r) -- (p);
      \end{tikzpicture}
      \raggedright
      \begin{pordentro}[Links verificados]
        \emph{App Links} (Android, arquivo \cod{assetlinks.json}) e \emph{Universal Links} (iOS, \cod{apple-app-site-association}) provam que o site e o app têm o mesmo dono.
      \end{pordentro}
    \end{column}
  \end{columns}
\end{frame}

\begin{frame}{Movimento: interpolação, curvas e molas}
  \begin{columns}[T]
    \begin{column}{0.44\textwidth}
      \centering
      \begin{tikzpicture}[x=3.4cm,y=1.9cm,baseline=(current bounding box.north)]
        \draw[->,cinzamedio] (0,0) -- (1.08,0) node[rotulo,right]{t};
        \draw[->,cinzamedio] (0,0) -- (0,1.35) node[rotulo,above]{progresso};
        \draw[cinzaclaro,dashed] (0,1) -- (1,1);
        \draw[cinzamedio,line width=0.8pt] plot[domain=0:1,samples=2] (\x,\x);
        \draw[azulpetroleo,line width=1pt] plot[domain=0:1,samples=40] (\x,{1-(1-\x)^3});
        \draw[ambar,line width=1pt] plot[domain=0:1,samples=80] (\x,{1-exp(-5*\x)*cos(deg(10*\x))});
        \node[rotulo,text=cinzamedio] at (0.82,0.62) {linear};
        \node[rotulo,text=azulpetroleo] at (0.3,0.88) {desacelera};
        \node[rotulo,text=ambar] at (0.42,1.3) {mola (passa e volta)};
      \end{tikzpicture}
      \raggedright
      \begin{definicao}[Interpolação]
        A cada quadro: \cod{valor = início + (fim - início)} $\times$ \cod{curva(t)}, com \cod{t} indo de 0 a 1 ao longo da animação.
      \end{definicao}
    \end{column}
    \begin{column}{0.54\textwidth}
      \begin{conceito}[Curvas e molas]
        Curvas de aceleração (\emph{easing}) imitam objetos reais: desaceleram ao chegar. \textbf{Molas} são definidas por rigidez e amortecimento, não por duração, e podem ser interrompidas sem saltos. O Material 3 Expressive (2025) adotou molas como base.
      \end{conceito}
      \begin{dica}[Movimento com propósito]
        Anime para orientar: de onde veio a tela, o que mudou. Transições de 150 a 400\,ms mantêm o app ágil.
      \end{dica}
      \begin{atencao}[Acessibilidade]
        Respeite a opção do sistema de remover ou reduzir animações: movimento excessivo causa desconforto em algumas pessoas.
      \end{atencao}
    \end{column}
  \end{columns}
\end{frame}


% =============================================================================
\section{Estado}

\begin{frame}{Estado efêmero e estado do app}
  \begin{columns}[T]
    \begin{column}{0.48\textwidth}
      \begin{conceito}[Estado efêmero (local)]
        Pertence a um único widget e pode ser perdido sem problema: aba selecionada, campo expandido, progresso de uma animação.\\[2pt]
        \textbf{Ferramenta}: \cod{StatefulWidget} + \cod{setState}.
      \end{conceito}
    \end{column}
    \begin{column}{0.48\textwidth}
      \begin{conceito}[Estado do app (compartilhado)]
        Usado por várias telas ou que precisa sobreviver: usuário logado, lista de bancadas, reservas, preferências.\\[2pt]
        \textbf{Ferramenta}: \cod{ChangeNotifier}, \emph{view models} e gerência de estado (Módulo 08).
      \end{conceito}
    \end{column}
  \end{columns}
  \vspace{4pt}
  \begin{center}
  \begin{tikzpicture}[node distance=10mm]
    \node[caixaambar,minimum width=2.6cm] (e) {\textbf{Estado}};
    \node[caixaescura,minimum width=2.6cm,right=of e] (b) {\cod{build()}};
    \node[caixa,minimum width=2.6cm,right=of b] (u) {\textbf{Tela}};
    \draw[seta] (e) -- (b); \draw[seta] (b) -- (u);
    \draw[setaambar,rounded corners] (u.south) -- ++(0,-5mm) -| node[rotulo,pos=0.25,below]{evento chama \cod{setState(\ldots)} e o \cod{build} roda de novo} (e.south);
  \end{tikzpicture}
  \end{center}
  \begin{dica}
    Comece com estado local e eleve-o só quando outra parte do app precisar dele.
  \end{dica}
\end{frame}

\begin{frame}{Ciclo de vida de um \cod{State}}
  \begin{columns}[c]
    \begin{column}{0.5\textwidth}
      \centering
      \begin{tikzpicture}[node distance=2mm,cb/.style={caixa,minimum width=3.6cm,font=\scriptsize}]
        \node[cb,caixaambar] (c) {\cod{createState()}};
        \node[cb,below=of c] (i) {\cod{initState()}};
        \node[cb,below=of i] (d) {\cod{didChangeDependencies()}};
        \node[cb,caixaescura,below=of d] (b) {\cod{build()}};
        \node[cb,below=of b] (u) {\cod{didUpdateWidget()} / \cod{setState()}};
        \node[cb,below=of u] (x) {\cod{dispose()}};
        \foreach \a/\bb in {c/i,i/d,d/b,u/x}{\draw[seta] (\a) -- (\bb);}
        \draw[seta] (b) -- (u);
        \draw[setaambar] (u.east) -- ++(4mm,0) |- (b.east);
      \end{tikzpicture}
    \end{column}
    \begin{column}{0.48\textwidth}
      \scriptsize
      \begin{tabularx}{\linewidth}{@{}l L@{}}
        \toprule
        \cabecalho{Método} & \cabecalho{Use para} \\
        \midrule
        \cod{initState} & criar controladores, iniciar inscrições (uma vez) \\
        \cod{didChangeDependencies} & reagir a \cod{Theme}, \cod{MediaQuery}, \cod{Provider} \\
        \cod{build} & descrever a tela (rápido, sem efeitos) \\
        \cod{didUpdateWidget} & o pai mudou os parâmetros \\
        \cod{dispose} & cancelar e liberar tudo \\
        \bottomrule
      \end{tabularx}
      \begin{atencao}
        \cod{build} pode rodar a cada quadro. Nada de requisições ou cálculos pesados dentro dele.
      \end{atencao}
    \end{column}
  \end{columns}
\end{frame}

\begin{frame}[fragile]{Elevando o estado}
  \begin{columns}[T]
    \begin{column}{0.52\textwidth}
\begin{codigo}[title=Componente sem estado]{Dart}
class CampoLimite extends StatelessWidget {
  const CampoLimite({super.key,
      required this.valor, required this.onChanged});
  final double valor;                   // estado desce
  final ValueChanged<double> onChanged; // evento sobe

  @override
  Widget build(BuildContext context) => Slider(
    value: valor, min: 20, max: 60, divisions: 40,
    label: '${valor.round()} °C',
    onChanged: onChanged,
  );
}

// o pai é o dono do estado
CampoLimite(
  valor: _limite,
  onChanged: (v) => setState(() => _limite = v),
);
\end{codigo}
    \end{column}
    \begin{column}{0.46\textwidth}
      \centering
      \begin{tikzpicture}[node distance=12mm]
        \node[caixaescura,minimum width=3.4cm] (p) {\cod{TelaAjustes}\\ \tiny dona do estado};
        \node[caixa,minimum width=3.4cm,below=of p] (f) {\cod{CampoLimite}\\ \tiny sem estado};
        \draw[seta] ([xshift=-8mm]p.south) -- node[rotulo,left]{valor} ([xshift=-8mm]f.north);
        \draw[setaambar] ([xshift=8mm]f.north) -- node[rotulo,right]{onChanged} ([xshift=8mm]p.south);
      \end{tikzpicture}
      \begin{conceito}[Fluxo unidirecional]
        O estado desce como parâmetro e os eventos sobem como funções. Há uma única fonte da verdade.
      \end{conceito}
    \end{column}
  \end{columns}
\end{frame}

\begin{frame}[fragile]{\cod{ValueNotifier} e \cod{ChangeNotifier}}
  \begin{columns}[T]
    \begin{column}{0.52\textwidth}
\begin{codigo}[title=Estado observável]{Dart}
class CarrinhoReservas extends ChangeNotifier {
  final _itens = <Reserva>[];
  List<Reserva> get itens => List.unmodifiable(_itens);

  void adicionar(Reserva r) {
    _itens.add(r);
    notifyListeners();          // avisa quem ouve
  }
}

final carrinho = CarrinhoReservas();

ListenableBuilder(
  listenable: carrinho,
  builder: (context, _) =>
      Badge(label: Text('${carrinho.itens.length}'),
            child: const Icon(Icons.event)),
);
\end{codigo}
    \end{column}
    \begin{column}{0.46\textwidth}
\begin{codigo}[title=Um único valor]{Dart}
final tema = ValueNotifier(ThemeMode.system);

ValueListenableBuilder(
  valueListenable: tema,
  builder: (context, modo, _) =>
      MaterialApp(themeMode: modo, ...),
);
tema.value = ThemeMode.dark;   // reconstrói
\end{codigo}
      \begin{conceito}
        Só o \cod{builder} é reconstruído quando o objeto notifica. É a base dos \emph{view models} recomendados pela documentação do Flutter.
      \end{conceito}
    \end{column}
  \end{columns}
\end{frame}

\begin{frame}[fragile]{Ciclo de vida do app}
  \begin{columns}[T]
    \begin{column}{0.54\textwidth}
\begin{codigo}[title=Pausando o sensor em segundo plano]{Dart}
class _MonitorState extends State<Monitor> {
  late final AppLifecycleListener _ouvinte;

  @override
  void initState() {
    super.initState();
    _ouvinte = AppLifecycleListener(
      onResume: () => sensor.retomar(),
      onPause: () => sensor.pausar(),
    );
  }

  @override
  void dispose() {
    _ouvinte.dispose();
    super.dispose();
  }
}
\end{codigo}
    \end{column}
    \begin{column}{0.44\textwidth}
      \scriptsize
      \begin{tabularx}{\linewidth}{@{}l L@{}}
        \toprule
        \cabecalho{Estado} & \cabecalho{Significado} \\
        \midrule
        \cod{resumed} & visível e recebendo toques \\
        \cod{inactive} & visível, sem foco (chamada, multitarefa) \\
        \cod{hidden} & não visível \\
        \cod{paused} & em segundo plano \\
        \cod{detached} & motor ativo sem tela \\
        \bottomrule
      \end{tabularx}
      \begin{atencao}[Morte do processo]
        Em segundo plano o sistema pode encerrar o app. Salve no banco ou no disco o que o usuário não pode perder (Módulo 10).
      \end{atencao}
    \end{column}
  \end{columns}
\end{frame}

% =============================================================================
\section{Navegação}

\begin{frame}[fragile]{\cod{Navigator}: a pilha de telas}
  \begin{columns}[T]
    \begin{column}{0.54\textwidth}
\begin{codigo}[title=Navegação imperativa]{Dart}
// abrir uma tela
Navigator.of(context).push(MaterialPageRoute(
  builder: (_) => TelaDetalhe(id: eq.patrimonio),
));

// voltar devolvendo um resultado
Navigator.of(context).pop(true);

// esperar o resultado
final salvou = await Navigator.of(context).push<bool>(
    MaterialPageRoute(builder: (_) => const TelaNova()));
if (salvou == true) recarregar();
\end{codigo}
      \begin{dica}
        Ótimo para fluxos simples. Para apps com URLs, \emph{deep links}, abas e login, use o \textbf{go\_router}.
      \end{dica}
    \end{column}
    \begin{column}{0.44\textwidth}
      \centering
      \begin{tikzpicture}[node distance=1mm]
        \node[caixa,minimum width=3cm] (a) {\cod{TelaLista}};
        \node[caixa,minimum width=3cm,above=of a] (b) {\cod{TelaDetalhe}};
        \node[caixaambarforte,minimum width=3cm,above=of b] (c) {\cod{TelaReserva}};
        \node[rotulo,right=2mm of c] {topo: visível};
        \node[pontilhado,fit=(a)(b)(c)] {};
        \node[rotulo,above=3mm of c] {\cod{push} empilha \quad \cod{pop} desempilha};
      \end{tikzpicture}
      \begin{conceito}
        O botão ou gesto voltar do sistema faz \cod{pop}. \cod{PopScope} permite interceptar (ex.: ``descartar alterações?'').
      \end{conceito}
    \end{column}
  \end{columns}
\end{frame}

\begin{frame}[fragile]{go\_router: rotas declarativas}
  \begin{columns}[T]
    \begin{column}{0.58\textwidth}
\begin{codigo}[title=lib/routing/router.dart]{Dart}
final router = GoRouter(
  initialLocation: '/painel',
  refreshListenable: auth,                // reavalia o login
  redirect: (context, state) =>
      auth.logado ? null : '/login',      // proteção
  routes: [
    GoRoute(path: '/login',
        builder: (_, __) => const TelaLogin()),
    GoRoute(path: '/painel',
        builder: (_, __) => const TelaPainel()),
    GoRoute(path: '/equipamentos',
      builder: (_, __) => const TelaEquipamentos(),
      routes: [
        GoRoute(path: ':id', builder: (_, state) =>
            TelaDetalhe(id: state.pathParameters['id']!)),
      ]),
  ],
);

MaterialApp.router(routerConfig: router);
\end{codigo}
    \end{column}
    \begin{column}{0.4\textwidth}
      \scriptsize
      \begin{tabularx}{\linewidth}{@{}l L@{}}
        \toprule
        \cabecalho{Chamada} & \cabecalho{Efeito} \\
        \midrule
        \cod{context.go('/painel')} & troca a pilha pela rota \\
        \cod{context.push('/x')} & empilha por cima \\
        \cod{context.pop()} & volta \\
        \cod{context.go('/equipamentos/42')} & caminho com parâmetro \\
        \bottomrule
      \end{tabularx}
      \begin{conceito}[Por que go\_router]
        Pacote mantido pela equipe do Flutter: URLs na web, \emph{deep links} (abrir o app num link), proteção de rotas e navegação por abas.
      \end{conceito}
    \end{column}
  \end{columns}
\end{frame}

\begin{frame}[fragile]{Abas com estado preservado}
  \begin{columns}[T]
    \begin{column}{0.58\textwidth}
\begin{codigo}[title=StatefulShellRoute]{Dart}
StatefulShellRoute.indexedStack(
  builder: (context, state, shell) => Scaffold(
    body: shell,
    bottomNavigationBar: NavigationBar(
      selectedIndex: shell.currentIndex,
      onDestinationSelected: shell.goBranch,
      destinations: const [
        NavigationDestination(
            icon: Icon(Icons.dashboard), label: 'Painel'),
        NavigationDestination(
            icon: Icon(Icons.memory), label: 'Equipamentos'),
        NavigationDestination(
            icon: Icon(Icons.settings), label: 'Ajustes'),
      ]),
  ),
  branches: [
    StatefulShellBranch(routes: [rotaPainel]),
    StatefulShellBranch(routes: [rotaEquipamentos]),
    StatefulShellBranch(routes: [rotaAjustes]),
  ],
);
\end{codigo}
    \end{column}
    \begin{column}{0.4\textwidth}
      \begin{conceito}[Uma pilha por aba]
        Cada \cod{StatefulShellBranch} guarda sua própria navegação: ao voltar para a aba Equipamentos, o usuário reencontra o detalhe que estava aberto.
      \end{conceito}
      \begin{dica}[Diálogos e avisos]
        \cod{showDialog}, \cod{showModalBottomSheet} e \cod{showSnackBar} (via \cod{ScaffoldMessenger}) também são camadas sobre a tela atual.
      \end{dica}
    \end{column}
  \end{columns}
\end{frame}

% =============================================================================
\section{Animações}

\begin{frame}[fragile]{Animações implícitas}
  \begin{columns}[T]
    \begin{column}{0.54\textwidth}
\begin{codigo}[title=Basta mudar o valor]{Dart}
AnimatedContainer(
  duration: const Duration(milliseconds: 300),
  curve: Curves.easeOut,
  height: expandido ? 160 : 72,
  color: emAlerta ? Colors.orange : Colors.white,
  child: conteudo,
);

AnimatedSwitcher(                   // troca com fade
  duration: const Duration(milliseconds: 250),
  child: Text('$temperatura °C',
      key: ValueKey(temperatura)),
);

TweenAnimationBuilder<double>(
  tween: Tween(begin: 0, end: nivel),
  duration: const Duration(seconds: 1),
  builder: (_, v, __) => LinearProgressIndicator(value: v),
);
\end{codigo}
    \end{column}
    \begin{column}{0.44\textwidth}
      \begin{conceito}[Implícita]
        Você muda uma propriedade e o widget anima sozinho do valor antigo ao novo. Cobre a maioria dos casos.
      \end{conceito}
      \scriptsize
      \begin{tabular}{@{}l l@{}}
        \toprule
        \cabecalho{Widget} & \cabecalho{Anima} \\
        \midrule
        \cod{AnimatedContainer} & tamanho, cor, borda \\
        \cod{AnimatedOpacity} & transparência \\
        \cod{AnimatedAlign} & posição \\
        \cod{AnimatedSwitcher} & troca de filho \\
        \cod{TweenAnimationBuilder} & qualquer valor \\
        \bottomrule
      \end{tabular}
    \end{column}
  \end{columns}
\end{frame}

\begin{frame}[fragile]{Animações explícitas e Hero}
  \begin{columns}[T]
    \begin{column}{0.56\textwidth}
\begin{codigo}[title=Indicador pulsando]{Dart}
class _AlertaState extends State<Alerta>
    with SingleTickerProviderStateMixin {
  late final _ctrl = AnimationController(vsync: this,
      duration: const Duration(seconds: 1))..repeat(reverse: true);
  late final _escala = Tween(begin: 1.0, end: 1.3)
      .animate(CurvedAnimation(parent: _ctrl,
          curve: Curves.easeInOut));

  @override
  void dispose() { _ctrl.dispose(); super.dispose(); }

  @override
  Widget build(BuildContext context) => ScaleTransition(
      scale: _escala,
      child: const Icon(Icons.warning, color: Colors.orange));
}
\end{codigo}
    \end{column}
    \begin{column}{0.42\textwidth}
      \begin{conceito}[Explícita]
        Um \cod{AnimationController} controla o tempo (iniciar, repetir, inverter). Use quando precisar de laços, sequências ou controle fino.
      \end{conceito}
      \begin{dica}[Hero]
        Envolva a foto do equipamento em \cod{Hero(tag: eq.patrimonio, \ldots)} na lista e no detalhe: ela ``voa'' entre as telas.
      \end{dica}
      \begin{dica}[Pacotes]
        \cod{animations}, \textbf{Rive} e \textbf{Lottie}.
      \end{dica}
    \end{column}
  \end{columns}
\end{frame}

% =============================================================================
\section{Fechamento}

\begin{frame}{Resumo}
  \begin{resumo}
    \begin{itemize}
      \item Estado efêmero com \cod{setState}; estado compartilhado com \cod{ChangeNotifier} e \emph{view models}.
      \item \cod{initState} cria, \cod{dispose} libera; \cod{build} só descreve a tela.
      \item Eleve o estado: valor desce, evento sobe, uma fonte da verdade.
      \item \cod{AppLifecycleListener} pausa recursos em segundo plano.
      \item go\_router: rotas declarativas, parâmetros, \cod{redirect} para login e abas com \cod{StatefulShellRoute}.
      \item Animações implícitas para o dia a dia; \cod{AnimationController} para controle fino; \cod{Hero} entre telas.
    \end{itemize}
  \end{resumo}
\end{frame}

\begin{frame}{Glossário (1/2)}
  \begin{glossario}
    Estado & valor que muda com o tempo e afeta o que a tela mostra \\
    Estado efêmero & pertence a um widget e pode se perder sem prejuízo \\
    Estado derivado & valor calculado a partir de outros; não se guarda \\
    Fonte única da verdade & cada dado tem um único dono que pode alterá-lo \\
    Fluxo unidirecional & o estado desce até a interface; os eventos sobem até o dono \\
    \cod{setState} & avisa que o estado mudou e agenda um novo \cod{build} \\
    \cod{initState}/\cod{dispose} & início e fim do ciclo de vida de um \cod{State} \\
    \cod{ChangeNotifier} & objeto que avisa seus ouvintes quando muda \\
    \cod{AppLifecycleState} & estado do app: \cod{resumed}, \cod{inactive}, \cod{hidden}, \cod{paused}, \cod{detached} \\
  \end{glossario}
\end{frame}

\begin{frame}{Glossário (2/2)}
  \begin{glossario}
    Morte do processo & encerramento do app em segundo plano para liberar memória \\
    \cod{RestorationMixin} & salva e devolve o estado depois da morte do processo \\
    Pilha de retorno & histórico LIFO dos destinos visitados \\
    go\_router & roteador declarativo baseado em caminhos de URL \\
    Link profundo & endereço que abre o app direto em um destino \\
    \cod{PopScope} & controla o voltar, inclusive o gesto preditivo do Android \\
    \cod{Curve} & função de aceleração de uma animação \\
    \cod{AnimationController} & controla tempo e direção de uma animação explícita \\
  \end{glossario}
\end{frame}

\begin{frame}{Questões de revisão}
  \begin{questoes}
    \item Por que um contador guardado em uma variável dentro de \cod{build} não funciona?
    \item Explique o fluxo unidirecional com um campo de busca e um \cod{ChangeNotifier}.
    \item Modele a tela de equipamentos como máquina de estados. Por que uma \cod{sealed class} e não três booleanos?
    \item O que acontece com o \cod{State} de um widget quando o Android encerra o processo em segundo plano? Como restaurar?
    \item Que recursos liberar quando o app passa para \cod{paused}?
    \item Desenhe a pilha depois de \cod{push} de Equipamentos e de Detalhe a partir do Painel, e de um \cod{pop}.
    \item Como o go\_router trata o link \cod{/reserva/42}? Que validações fazer?
    \item Por que preferir uma mola a uma curva de duração fixa quando o usuário pode interromper o gesto?
  \end{questoes}
\end{frame}

\begin{frame}{Exercícios}
  \begin{exercicio}[Exercício 1 --- Ciclo de vida]
    Registre um \cod{debugPrint} em cada método do ciclo de vida de um \cod{State}. Navegue, gire a tela e volte; descreva a sequência.
  \end{exercicio}
  \begin{exercicio}[Exercício 2 --- Navegação]
    Configure o go\_router com login, três abas e a rota \cod{/equipamentos/:id}. Abra o app pelo navegador já na URL de um equipamento.
  \end{exercicio}
  \begin{exercicio}[Exercício 3 --- Animação]
    Faça o cartão da bancada mudar de cor com \cod{AnimatedContainer} quando a temperatura passar do limite e pulsar um ícone de alerta.
  \end{exercicio}
  \begin{projeto}[LabMonitor --- Etapa 7]
    Implemente a navegação completa: login, Painel, Equipamentos (lista $\rightarrow$ detalhe $\rightarrow$ reserva) e Ajustes, com abas preservadas e animação de alerta.
  \end{projeto}
\end{frame}

\begin{frame}{Para aprofundar}
  \small
  \begin{itemize}
    \item Estado efêmero e de app: \url{https://docs.flutter.dev/data-and-backend/state-mgmt/ephemeral-vs-app}
    \item Navegação: \url{https://docs.flutter.dev/ui/navigation}
    \item go\_router: \url{https://pub.dev/packages/go_router}
    \item Animações: \url{https://docs.flutter.dev/ui/animations}
    \item Ciclo de vida do app: \url{https://api.flutter.dev/flutter/widgets/AppLifecycleListener-class.html}
  \end{itemize}
\end{frame}

\begin{frame}{Respostas das questões de revisão}
  \begin{respostas}
    \item \cod{build} roda de novo a cada reconstrução e a variável volta ao valor inicial. O estado deve ficar em um campo do \cod{State}, alterado com \cod{setState}, ou em um \cod{ChangeNotifier}.
    \item O campo mostra o texto vindo do \emph{notifier} (estado que desce) e chama \cod{buscar(texto)} ao digitar (evento que sobe); o \emph{notifier} filtra a lista, chama \cod{notifyListeners} e a tela se reconstrói.
    \item Estados \cod{Carregando}, \cod{Sucesso(itens)}, \cod{Vazio} e \cod{Erro(msg)}. A classe selada só admite esses quatro e o \cod{switch} obriga a tratar todos; booleanos permitem combinações impossíveis.
    \item Ele se perde: o processo recomeça do zero. Com \cod{RestorationMixin} e \cod{restorationScopeId}, os valores registrados são salvos e devolvidos; dados importantes também podem ir para o disco.
    \item Câmera, sensores, fluxos de localização e animações contínuas; também é hora de salvar rascunhos. Tudo é retomado em \cod{resumed}.
    \item Depois dos \cod{push}: Painel, Equipamentos, Detalhe (topo). Depois do \cod{pop}: Painel, Equipamentos.
    \item A rota \cod{/reserva/:id} recebe \cod{42} em \cod{state.pathParameters}, e rotas aninhadas deixam o Painel abaixo na pilha. Validar se a reserva existe e se o usuário pode vê-la, com \cod{redirect} para o login se preciso.
    \item A mola parte da posição e da velocidade atuais e segue a física até o novo alvo; a curva recomeça com duração fixa e muda a velocidade de forma brusca.
  \end{respostas}
\end{frame}

\encerramento{Módulo 08 --- Arquitetura e gerência de estado: MVVM, repositórios, Provider, Riverpod e Bloc.}

\end{document}
